\newcommand{\objone}{
The objective of this practical is to implement the Gauss--Seidel iterative method 
to solve a system of linear equations $\mathbf{A}\mathbf{x}=\mathbf{b}$. 
The method updates each component of the solution vector successively using the formula
\[
x_i^{(k+1)} = \frac{1}{a_{ii}} \left( b_i - \sum_{j<i} a_{ij}x_j^{(k+1)} - \sum_{j>i} a_{ij}x_j^{(k)} \right), 
\quad i=1,2,\dots,n.
\]
The aim is to develop a program for this algorithm, apply it to suitable test 
systems, analyze its convergence behavior, and compare the number of iterations 
and accuracy with other iterative methods such as the Jacobi method.
}


\newcommand{\objtwo}{
The objective of this practical is to write and implement a program to find the root 
of a nonlinear equation $f(x)=0$ using the Bisection Method. The method repeatedly 
divides the interval $[a,b]$ into two halves and selects the subinterval where a 
sign change occurs, based on the formula
\[
c = \frac{a+b}{2},
\]
where $c$ is the midpoint of the current interval. The process is continued until 
the root is obtained within a desired tolerance. This practical also aims to study 
the convergence, efficiency, and limitations of the Bisection Method compared to 
other root-finding techniques.
}


\newcommand{\objthree}{
The objective of this practical is to implement the Newton--Raphson method to find the 
root of a given nonlinear equation $f(x)=0$. The method uses the tangent line at an 
approximate root to iteratively improve the solution, based on the formula
\[
x_{n+1} = x_n - \frac{f(x_n)}{f'(x_n)}.
\]
The aim is to develop a program for this algorithm, apply it to suitable test 
functions, and analyze its convergence properties. This practical also focuses 
on studying the rapid convergence of the Newton--Raphson method and comparing 
its efficiency with bracketing methods such as the Bisection Method.
}


\newcommand{\objfour}{
The objective of this practical is to develop a program to find the inverse of a 
matrix using the Gauss--Jordan elimination method. The method is based on forming 
the augmented matrix $[A|I]$ and applying row operations to transform it into 
$[I|A^{-1}]$, where $A^{-1}$ is the required inverse matrix.  

Mathematically,
\[
[A|I] \;\;\xrightarrow{\text{row operations}}\;\; [I|A^{-1}].
\]

The aim is to implement this algorithm, practice systematic row transformations, 
and verify the correctness of the inverse by checking that $A A^{-1} = I$. This 
practical also highlights the importance of matrix inversion in solving systems 
of linear equations and other computational applications.
}


\newcommand{\objfive}{
The objective of this practical is to implement the Lagrange interpolation method 
to approximate a function from a given set of data points $(x_0, y_0), (x_1, y_1), 
\dots, (x_n, y_n)$. The method constructs a polynomial of degree $n$ that passes 
through all the given points, using the formula
\[
P(x) = \sum_{i=0}^{n} y_i \, L_i(x),
\quad \text{where} \quad
L_i(x) = \prod_{\substack{j=0 \\ j \neq i}}^{n} \frac{x - x_j}{x_i - x_j}.
\]

The aim is to develop a program to compute the interpolating polynomial, 
evaluate it at desired points, and compare the interpolated values with the 
original function. This practical also emphasizes the application of interpolation 
in numerical analysis and its importance in approximating unknown values from 
discrete data.
}

\newcommand{\objsix}{
The objective of this practical is to write programs to approximate the value of a 
definite integral $\int_a^b f(x)\,dx$ using the Trapezoidal Rule and Simpson's 
$\tfrac{1}{3}$ Rule.  

The Trapezoidal Rule is based on approximating the area under the curve with 
trapezoids, given by
\[
\int_a^b f(x)\,dx \approx \frac{h}{2} \left[ f(x_0) + 2\sum_{i=1}^{n-1} f(x_i) + f(x_n) \right],
\]
where $h = \frac{b-a}{n}$ and $n$ is the number of subintervals.  

Simpson's $\tfrac{1}{3}$ Rule uses parabolic arcs for approximation, expressed as
\[
\int_a^b f(x)\,dx \approx \frac{h}{3} \left[ f(x_0) + 4\sum_{\substack{i=1 \\ i\;\text{odd}}}^{n-1} f(x_i) 
+ 2\sum_{\substack{i=2 \\ i\;\text{even}}}^{n-2} f(x_i) + f(x_n) \right],
\]
where $n$ is even.  

The aim is to implement these rules in a program, compare the approximations with the 
exact values, analyze their accuracy, and study the effect of increasing the number of 
subintervals on the error.
}

\newcommand{\objseven}{
The objective of this practical is to implement numerical methods for solving 
ordinary differential equations (ODEs) of the form
\[
\frac{dy}{dx} = f(x,y), \quad y(x_0) = y_0,
\]
using Euler's Method and Runge--Kutta Methods of second and fourth order.  

Euler's Method is given by
\[
y_{n+1} = y_n + h f(x_n, y_n),
\]
where $h$ is the step size.  

The Second-Order Runge--Kutta Method (RK2) can be expressed as
\[
k_1 = h f(x_n, y_n), \quad
k_2 = h f(x_n + h, \, y_n + k_1), \quad
y_{n+1} = y_n + \tfrac{1}{2}(k_1 + k_2).
\]

The Fourth-Order Runge--Kutta Method (RK4) is defined as
\[
\begin{aligned}
k_1 &= h f(x_n, y_n), \\
k_2 &= h f(x_n + \tfrac{h}{2}, \, y_n + \tfrac{k_1}{2}), \\
k_3 &= h f(x_n + \tfrac{h}{2}, \, y_n + \tfrac{k_2}{2}), \\
k_4 &= h f(x_n + h, \, y_n + k_3), \\
y_{n+1} &= y_n + \tfrac{1}{6}(k_1 + 2k_2 + 2k_3 + k_4).
\end{aligned}
\]

The aim is to implement these algorithms, compare their numerical solutions with 
analytical solutions where available, study their accuracy and stability, and 
observe the effect of step size on the error.
}

\newcommand{\objeight}{
The objective of this practical is to fit linear, exponential, and polynomial curves 
to a given set of discrete data points using the Least Squares Method. The method 
minimizes the sum of squared errors between the observed data and the fitted curve.  

For Linear Curve Fitting ($y = a + bx$), the normal equations are
\[
\sum y = na + b \sum x, 
\qquad 
\sum xy = a \sum x + b \sum x^2.
\]

For Exponential Curve Fitting ($y = ae^{bx}$), taking logarithms gives
\[
\ln y = \ln a + bx,
\]
which reduces to a linear least squares problem.  

For Polynomial Curve Fitting ($y = a_0 + a_1x + a_2x^2 + \dots + a_mx^m$), the normal 
equations are obtained as
\[
\begin{align}
\sum x^k y &= a_0 \sum x^k + a_1 \sum x^{k+1} + \dots + a_m \sum x^{k+m}, \\
           &\quad k = 0,1,\dots,m.
\end{align}
\]

The aim is to implement programs for these curve fitting techniques, determine the 
best-fit parameters, evaluate the goodness of fit, and compare fitted curves with the 
original data graphically.
}

\newcommand{\objnine}{
The objective of this practical is to solve a system of linear equations 
$\mathbf{A}\mathbf{x}=\mathbf{b}$ using the Gauss Elimination Method and the 
Gauss--Jordan Method. These are direct methods based on systematic row operations 
to reduce the coefficient matrix to a simpler form.  

In the Gauss Elimination Method, the system is transformed into an upper triangular 
system by forward elimination, followed by back substitution:
\[
U\mathbf{x} = \mathbf{c},
\]
where $U$ is upper triangular.  

In the Gauss--Jordan Method, the augmented matrix $[A|b]$ is reduced by row 
operations into the form
\[
[I|x],
\]
where $I$ is the identity matrix and $\mathbf{x}$ is the solution vector.  

The aim is to implement programs for both methods, apply them to test systems, 
and compare their efficiency, computational steps, and accuracy.
}